\documentclass{article}
\usepackage{graphicx}
\usepackage[a4paper, left=2cm, right=2cm, top=2.5cm, bottom=2.5cm]{geometry}

\usepackage[utf8]{inputenc}
\usepackage[greek,english]{babel}
\usepackage{alphabeta}
\usepackage{amsmath}
\usepackage{caption}
\usepackage{subcaption}
\usepackage{xcolor}
\usepackage{mathabx}
\usepackage{amssymb}
\usepackage{enumitem}
\usepackage{float}

\captionsetup{labelformat=empty}
\subcaptionsetup{labelformat=empty}

\title{\textbf{Προσαρμοστικός Έλεγχος - Εργαστήριο:\\Coupled DC Motors}}
\author{\small{\textbf{Μέλη Ομάδας}:} \\ \small{Βραχωρίτη Αλεξάνδρα (A.M.: 1092793)} \\ \small{Μπαθρέλλος Μιχαήλ (Α.Μ.: 1092665)} \\ \small{Σταναθιώτη Παναγιώτα (Α.Μ.: 1092673)}}
\date{\small{10 Φεβρουαρίου 2026}}

\begin{document}
\maketitle 
    \section{Αναγνώριση και ανάλυση του συστήματος}
    Στόχος αυτού του εργαστηρίου είναι η προσέγγιση του συστήματος με μια (κανονικοποιημένη) συνάρτηση μεταφοράς πρώτης τάξης:
    
    \begin{equation*}
        G(s) = \frac{K}{τs+1}
    \end{equation*}

    \hfill\break
    διαλέγοντας ως ονομαστική λειτουργία: 
    \begin{itemize}
        \item \(PWM_{nom} = 200\)
        \item \(RPM_{nom} = 19.5\)
        \item \(error_{nom} = 2\)
        \item \({\Delta\Omega}_{max} = 2\)
        \item \({\Delta u}_{max} = 55\)
    \end{itemize}
    
    \hfill\break
    Θα χρειαστεί, λοιπόν, να υπολογίσουμε το κέρδος και την χρονική σταθερά του συστήματος, επομένως πρέπει να δοκιμάσουμε διάφορες τιμές PWM (θετικές και αρνητικές) στο σύστημα και να καταγράψουμε τις στροφές (rpm) του κινητήρα, αλλά και τον χρόνο της απόκρισης από το 0\% μέχρι το 63\% της μέγιστης τιμής της εξόδου. Παρακάτω παρουσιάζονται οι πίνακες των μετρήσεων, ενώ οι μετρήσεις πάρθηκαν για εξωτερική τροφοδοσία από \(2.75\,V\) έως \(5\,V\), με την αντίσταση του ποτενσιόμετρου να είναι σταθερή στα \(R = 5\,k\Omega\).

    \begin{table}[!h]
        \centering
        \begin{tabular}{|c|c|c|c|c|}
        \hline
            \(V_{in}\) (Volt) & \(PWM_{in}\) & \(\Omega_{out}\) (rpm) & \(\tau\) (sec) & \(Κ = \frac{\Delta\Omega/\Delta\Omega_{max}}{\Delta u/{\Delta u}_{max}}\)\\ \hline
            2.75 & 140 &  5.4 & 0.3 & 6.46\\ \hline
            2.94 & 150 &  8.5 & 0.4 & 6.05\\ \hline
            3.14 & 160 & 11.1 & 0.4 & 5.78\\ \hline
            3.33 & 170 & 11.7 & 0.4 & 7.15\\ \hline
            3.53 & 180 & 15.4 & 0.3 & 5.64\\ \hline
            3.73 & 190 & 16.8 & 0.4 & 7.43\\ \hline
            3.92 & 200 & 19.5 &   0 & -   \\ \hline
            4.12 & 210 &   20 & 0.4 & 1.38\\ \hline
            4.31 & 220 & 21.4 & 0.5 & 2.61\\ \hline
            4.51 & 230 & 22.7 & 0.5 & 2.93\\ \hline
            4.71 & 240 &   24 & 0.4 & 3.09\\ \hline
            5.00 & 255 & 27.4 & 0.3 & 3.95\\ \hline
        \end{tabular}
        \caption{\textbf{Πίνακας 1.} Εδώ επιλέξαμε nominal PWM \(= 200\), δίνοντας στο σύστημά μας μόνο θετικές τιμές PWM.}
    \end{table}

    \begin{table}[!h]
        \centering
        \begin{tabular}{|c|c|c|c|c|}
        \hline
            \(V_{in}\) (Volt) & \(PWM_{in}\) & \(\Omega_{out}\) (rpm) & \(\tau\) (sec) & \(Κ = \frac{\Delta\Omega/\Delta\Omega_{max}}{\Delta u/{\Delta u}_{max}}\)\\ \hline
            -2.75 & -140 &  6.5 & 0.5 & 5.27\\ \hline
            -2.94 & -150 &  9.4 & 0.3 & 4.73\\ \hline
            -3.14 & -160 & 11.4 & 0.3 & 4.54\\ \hline
            -3.33 & -170 & 13.4 & 0.3 & 4.22\\ \hline
            -3.53 & -180 & 14.8 & 0.3 & 4.40\\ \hline
            -3.73 & -190 & 16.3 & 0.4 & 4.68\\ \hline
            -3.92 & -200 &   18 &   0 & -   \\ \hline
            -4.12 & -210 & 19.7 & 0.5 & 4.68\\ \hline
            -4.31 & -220 & 20.6 & 0.5 & 3.58\\ \hline
            -4.51 & -230 & 21.4 & 0.4 & 3.12\\ \hline
            -4.71 & -240 & 22.6 & 0.4 & 3.16\\ \hline
            -5.00 & -255 &   26 & 0.3 & 4.00
        \end{tabular}
        \caption{\textbf{Πίνακας 2.} Εδώ επιλέξαμε nominal PWM \(= -200\), δίνοντας στο σύστημά μας μόνο αρνητικές τιμές PWM.}
    \end{table}

    \hfill\break
    Τελικά υπολογίστηκε: \(K = 4.49\), \(\tau = 0.39 \Rightarrow pole = -2.59\). Έτσι, η κανονικοποιημένη συνάρτηση μεταφοράς που προσεγγίζει το σύστημα είναι:
    
    \begin{equation*}
        \boxed{G(s) = \frac{11.513}{s + 2.59}}
    \end{equation*}
    
    \noindent
    ή αλλιώς:
    
    \begin{align*}
        \dot{x} = -2.59x + 11.513u
    \end{align*}
    
    \noindent
    Η μη κανονικοποιημένη συνάρτηση μεταφοράς είναι αντίστοιχα:
    \begin{equation*}
        \boxed{G(s) = \frac{0.418}{s + 2.59}}
    \end{equation*}

    \noindent
    ή αλλιώς

    \noindent
    \begin{align*}
        \dot{x} = -2.59x + 0.418u
    \end{align*}
    
    % \hfill\break
    % Παρακάτω θα δούμε την απόκριση του συστήματος για κάποιες εισόδους, ώστε να επιβεβαιώσουμε ότι το αποτέλεσμά μας είναι ορθό.

    % EIKONES

    % \hfill\break
    % Βάσει των παραπάνω, παρατηρούμε πως το σύστημα λειτουργεί σαν βαθυπερατό φίλτρο, επιβεβαιώνοντας την ανάλυσή μας.

    % ----------------------------------------------------- PI ----------------------------------------------------- %
    \section{Εργαστήριο I.1:  Έλεγχος ταχύτητας με PI ελεγκτή}
    Παρακάτω παρουσιάζονται τα αποτελέσματα προσομοίωσης για την παρακολούθηση της βηματικής συνάρτησης με εναλλαγές από τα 20 στα 15 rpm κ.ο.κ.

    \begin{figure}[h!]
        \centering
        \begin{subfigure}{0.45\textwidth}
            \centering
            \includegraphics[width=\textwidth]{1-P_controller.png}
            \caption{\textbf{Διάγραμμα 2.1.1} Αναλογικός ελεγκτής \(P\) με \(K_p = 10\).}
        \end{subfigure}
        \,\,\,\,\,\,
        \begin{subfigure}{0.45\textwidth}
            \centering
            \includegraphics[width=\textwidth]{2-PI_controller.png}
            \caption{\textbf{Διάγραμμα 2.1.2} Αναλογικός - Ολοκληρωτικός ελεγκτής PI με \(K_i = 3\).}
        \end{subfigure}
    \end{figure}
    
    \hfill\break
    Με τον αναλογικό ελεγκτή (P) παρατηρούμε την εμφάνιση του μόνιμου σφάλματος, το οποίο διορθώνεται με την πρόσθεση του \(I\) όρου. Έτσι, ο PI ελεγκτής απαλείφει το σφάλμα μόνιμης κατάστασης και άρα μπορούμε να ελέγξουμε ικανοποιητικά το σύστημα. Παρακάτω θα εξετάσουμε τις δυσλειτουργίες του PI ελεγκτή.
    
    \subsection{Επίδραση του ποτενσιόμετρου στη δυναμική του συστήματος}
    Η συνάρτηση μεταφοράς του συστήματος υπολογίστηκε για σταθερή αντίσταση του ποτενσιόμετρου \(R = 5k\Omega\), με το tuning του PI ελεγκτή να είναι προσαρμοσμένο σε αυτό το σύστημα. Μεταβάλλοντας την αντίσταση, το σύστημα θα αλλάξει και ο PI ελεγκτής δεν δουλεύει πια. Αυτό φαίνεται στα παρακάτω διαγράμματα:

    \begin{figure}[h!]
        \centering
        \includegraphics[width=0.64\textwidth]{3-PI_controller_poten.png}
        \caption{\textbf{Διάγραμμα 2.1.3} Αποτέλεσμα προσομοίωσης με αλλαγή της αντίστασης μέσω ποτενσιόμετρου.}
    \end{figure}

    \noindent
    Η αντίσταση του ποτενσιόμετρου προσομοιάζει παραμετρική αβεβαιότητα. Στα επόμενα εργαστήρια θα μελετήσουμε προσαρμοστικούς ελεγκτές οι οποίοι αντιμετωπίζουν αυτού του είδους αβεβαιότητες.

    % ----------------------------------------------------- MIT ----------------------------------------------------- %
    \section{Εργαστήριο I.2: Έλεγχος ταχύτητας με τη μέθοδο του κανόνα MIT}
    Ο κανόνας MIT είναι ένας από τους πρώτους νόμους προσαρμογής που χρησιμοποιήθηκαν στον προσαρμοστικό έλεγχο και βασίζεται στην ιδέα της ελαχιστοποίησης του κριτηρίου $J(\theta) = \frac{1}{2}e_0^2(\theta)$, όπου $e_0^2(\theta)$ το σφάλμα παρακολούθησης. Η προσαρμογή των αγνώστων παραμέτρων $\theta$ προκύπτει με χρήση gradient descent μεθόδου, γεγονός που προσδίδει στον κανόνα MIT 
    την κλειστή μαθηματική έκφραση: 
    
    \begin{align*}
       \frac{d\theta}{dt} = -\gamma e_0(\theta) \frac{\partial}{\partial \theta} y_p(\theta).
    \end{align*}

    \noindent
    Ωστόσο, στον κανόνα MIT είναι δύσκολη η απόδειξη global σύγκλισης ή ευστάθειας, 
    αφού δεν στηρίζεται στη θεωρία Lyapunov. Η μεταβολή των παραμέτρων εξαρτάται 
    από την επιλογή του κέρδους προσαρμογής, του σήματος αναφοράς και 
    της δυναμικής του plant.\\

    \noindent
    Το σύστημα αναφοράς όπου θα ακολουθεί το σύστημα μέσω του ελεγκτή είναι:

    \begin{align*}
        \dot{x}_m &= a_mx + b_mr \\
        &=-0.9x + 0.9r
    \end{align*}
    Ο νόμος ελέγχου είναι:
    \begin{align*}
        u = \hat{θ}_1r - \hat{θ}_2x 
    \end{align*}
    και ο νόμος προσαρμογής είναι:
    \begin{align*}
        \frac{d\theta}{dt} = -\gamma e_0(\theta) \frac{\partial}{\partial \theta} e_0(\theta).
    \end{align*}
    
    % \subsection{Παρακολούθηση βηματικής συνάρτησης}
    \subsection{Σήμα αναφοράς: βηματική συνάρτηση}
    \begin{figure}[H]
        \centering
        \begin{subfigure}{0.45\textwidth}
            \centering
            \includegraphics[width=\textwidth]{1-mit_g05.png}
            \caption{\textbf{Διάγραμμα 3.1.1} Παρακολούθηση του σήματος αναφοράς \(r(t) = 16\).}
        \end{subfigure}
        \,\,\,\,\,\,
        \begin{subfigure}{0.45\textwidth}
            \centering
            \includegraphics[width=\textwidth]{2-mit_g10.png}
            \caption{\textbf{Διάγραμμα 3.1.2} Παρακολούθηση διάφορων βηματικών συναρτήσεων.}
        \end{subfigure}
    \end{figure}
    \noindent
    Παρατηρείται πως ο ελεγκτής παρακολουθεί ικανοποιητικά βηματικές εισόδους. Επίσης, οι παράμετροι συγκλίνουν σε σταθερές τιμές για επιλεγμένα κέρδη προσαρμογής \(γ = 5, γ = 10\).

    \begin{figure}[H]
        \centering
        \includegraphics[width=0.64\linewidth]{8-mit_g1_a0001.png}
        \caption{\textbf{Διάγραμμα 3.1.3} Υλοποίηση κανονικοποιημένου MIT κανόνα και παρακολούθηση διαφόρων βηματικών συναρτήσεων}
        \label{fig:placeholder}
    \end{figure}
    
    \noindent
    Η υλοποίηση του κανονικοποιημένου MIT κανόνα αποσκοπεί στη μείωση της ευαισθησίας του νόμου προσαρμογής στο πλάτος των σημάτων αναφοράς. Με την κανονικοποίηση αποφεύγονται υπερβολικά μεγάλες μεταβολές των παραμέτρων, βελτιώνοντας την ευστάθεια του προσαρμοστικού ελεγκτή.

    % \subsection{Παρακολούθηση ημιτονοειδούς συνάρτησης}
    \subsection{Σήμα αναφοράς: ημιτονοειδής συνάρτηση}
    \begin{figure}[H]
        \centering
        \begin{subfigure}{0.45\textwidth}
            \centering
            \includegraphics[width=\textwidth]{13-mit_16plus3sin04_g1.png}
            \caption{\textbf{Διάγραμμα 3.2.1} Παρακολούθηση του σήματος αναφοράς \(r(t) = 16 + 3sin(0.4t)\).}
        \end{subfigure}
        \,\,\,\,\,\,
        \begin{subfigure}{0.45\textwidth}
            \centering
            \includegraphics[width=\textwidth]{11-mit_g1_sin3.png}
            \caption{\textbf{Διάγραμμα 3.2.2} Παρακολούθηση του σήματος αναφοράς \(r(t) = 16 + sin(3t)\).}
        \end{subfigure}
    \end{figure}
    \noindent
    Ο ελεγκτής καταφέρνει να παρακολουθήσει ημιτονοειδή σήματα αναφοράς. Παρατηρείται πως όσο μικρότερη είναι η συχνότητα, τόσο καλύτερη η παρακολούθηση. Αυτό οφείλεται στην αδυναμία του plant να παρακολουθήσει ένα γρήγορα μεταβαλλόμενο μοντέλο αναφοράς. 

    \subsection{Συνθήκες ευστάθειας}
    \begin{figure}[H]
        \centering
        \begin{subfigure}{0.45\textwidth}
            \centering
            \includegraphics[width=\textwidth]{5-mit_g500.png}
            \caption{\textbf{Διάγραμμα 3.3.1} Παρακολούθηση του σήματος αναφοράς \(r(t) = 20\) για \(γ=500\).}
        \end{subfigure}
        \,\,\,\,\,\,
        \begin{subfigure}{0.45\textwidth}
            \centering
            \includegraphics[width=\textwidth]{12-mit_g1_sin10.png}
            \caption{\textbf{Διάγραμμα 3.3.2} Παρακολούθηση του σήματος αναφοράς \(r(t) = 16 + sin(10t)\) για \(γ=1\).}
        \end{subfigure}
    \end{figure}
    
    Για τον χαρακτηρισμό της ευστάθειας του συστήματος, μελετήθηκε η επίδραση δύο παραμέτρων: του προσαρμοστικού κέρδους $\gamma$ και της συχνότητας του σήματος αναφοράς $\omega$. Στην περίπτωση του μη κανονικοποιημένου MIT κανόνα, το πλάτος $A$ του σήματος αναφοράς αποτελεί επίσης κρίσιμη παράμετρο, καθώς επηρεάζει τον ρυθμό προσαρμογής και μπορεί να οδηγήσει σε αστάθεια.
    
    \hfill\break
    Παρατηρήθηκε ότι για τιμές $\gamma > 500$ οι παράμετροι αυξάνονται ανεξέλεγκτα, με αποτέλεσμα την απώλεια ευστάθειας του συστήματος. Επιπλέον, για συχνότητες $\omega > 10$ εμφανίζεται σημαντικό σφάλμα στη μόνιμη κατάσταση. Στον κανονικοποιημένο MIT κανόνα, η επίδραση του πλάτους $A$ στον νόμο προσαρμογής περιορίζεται σημαντικά. Όμως, για πολύ μεγάλα πλάτη το σύστημα δεν θα μπορεί να ακολουθήσει το σήμα αναφοράς λόγω φυσικών περιορισμών του ενεργοποιητή (ανεπαρκής διαθέσιμη ροπή του κινητήρα).

    
    % ----------------------------------------------------- MRAC ----------------------------------------------------- %
    \section{Εργαστήριο I.3: Έλεγχος ταχύτητας με τη μέθοδο MRAC (βαθμωτή περίπτωση)}
    Η μέθοδος MRAC είναι direct μέθοδος προσαρμοστικού ελέγχου καθώς υπολογίζει κατευθείαν τα κέρδη του ελεγκτή και βασίζεται στην θεωρία ευστάθειας Lyapunov. Για την εφαρμογή αυτού του ελεγκτή αναγκαία είναι η γνώση του προσήμου του \(b\). Το σύστημα αναφοράς όπου θα ακολουθεί το σύστημα μέσω του ελεγκτή είναι:
    \begin{equation*}
        \dot{x}_m = a_mx + b_mr =-0.9x + 0.9r
    \end{equation*}

    \noindent
    Ο νόμος ελέγχου του MRAC είναι:
    
    \begin{equation*}
        u = \hat{k}_xx + \hat{k}_rr
    \end{equation*}
    
    \noindent
    και οι νόμοι προσαρμογής είναι:
    
    \begin{equation*}
        \begin{aligned}
            \dot{\hat{k}}_x &= -γ_x\,x\,e\,sgn(b) \\
            \dot{\hat{k}}_r &= -γ_r\,r\,e\,sgn(b)
        \end{aligned}
    \end{equation*}
    
    % \subsection{Παρακολούθηση βηματικής συνάρτησης}
    \subsection{Σήμα αναφοράς: βηματική συνάρτηση}
    \begin{figure}[H]
        \centering
        \begin{subfigure}{0.45\textwidth}
            \centering
            \includegraphics[width=\textwidth]{1-mrac_one_step_follow.png}
            \caption{\textbf{Διάγραμμα 4.1.1} Παρακολούθηση του σήματος αναφοράς \(r(t) = 18\).}
        \end{subfigure}
        \,\,\,\,\,\,
        \begin{subfigure}{0.45\textwidth}
            \centering
            \includegraphics[width=\textwidth]{2-mrac_follow_step_g001.png}
            \caption{\textbf{Διάγραμμα 4.1.2} Παρακολούθηση διάφορων βηματικών συναρτήσεων.}
        \end{subfigure}
    \end{figure}

    \begin{figure}[H]
        \centering
        \includegraphics[width=0.59\linewidth]{3-mrac_potensiometer.png}
        \caption{\textbf{Διάγραμμα 4.1.3} Αποδοτικότητα ελεγκτή MRAC όταν αλλάζουν οι παράμετροι του συστήματος με αλλαγή αντίστασης ποτενσιόμετρου.}
    \end{figure}

    % \noindent
    % Παρατηρείται πως ο ελεγκτής παρακολουθεί ικανοποιητικά βηματικές εισόδους. Επίσης, οι παράμετροι έχουν συγκλίνει σε σταθερές τιμές, όμως λόγω της έλλειψης συχνοτικού περιεχόμενου δεν υπάρχει εγγύηση σύκγλισης στις ιδανικές τιμές.  ---> το έβαλα μετά τις εικόνες για να βγαίνει καλά η στοίχιση, οπότε γι' αυτό το άλλαξα λίγο
    % \noindent
    % Ο προσαρμοστικός ελεγκτής δουλεύει παρά την αλλαγή των παραμέτρων του συστήματος, επιβεβαιώνοντας την θεωρία.

    \noindent
    Από τα Διαγράμματα 4.1.1 και 4.1.2 παρατηρείται πως ο ελεγκτής παρακολουθεί ικανοποιητικά βηματικές εισόδους. Επίσης, οι παράμετροι έχουν συγκλίνει σε σταθερές τιμές, όμως λόγω της έλλειψης συχνοτικού περιεχόμενου δεν υπάρχει εγγύηση σύγκλισης στις ιδανικές τιμές.
    \noindent
    Τέλος, από το Διάγραμμα 4.1.3 είναι εμφανές πώς ο προσαρμοστικός ελεγκτής δουλεύει παρά την αλλαγή των παραμέτρων του συστήματος, επιβεβαιώνοντας την θεωρία. 

    % \subsection{Παρακολούθηση Ημιτονοειδούς συνάρτησης}
    \subsection{Σήμα αναφοράς: ημιτονοειδής συνάρτηση}
    \begin{figure}[H]
        \centering
        \begin{subfigure}{0.45\textwidth}
            \centering
            \includegraphics[width=\textwidth]{4-mrac_follow_3sin2t_g001.png}
            \caption{\textbf{Διάγραμμα 4.2.1} Παρακολούθηση του σήματος αναφοράς \(r(t) = 15 + 3sin(2t)\).}
        \end{subfigure}
        \,\,\,\,\,\,
        \begin{subfigure}{0.45\textwidth}
            \centering
            \includegraphics[width=\textwidth]{8-mrac_PE_sin_am2_bm2.png}
            \caption{\textbf{Διάγραμμα 4.2.2} Παρακολούθηση του σήματος αναφοράς \(r(t) = 17 + 3sin(0.7t)\).}
        \end{subfigure}
    \end{figure}
    \noindent
    Ο ελεγκτής καταφέρνει να παρακολουθήσει ημιτονοειδή σήματα αναφοράς. Παρατηρείται πως όσο μικρότερη είναι η συχνότητα, τόσο καλύτερη η παρακολούθηση. Στην περίπτωση αυτή, το σήμα αναφοράς έχει μία συχνότητα, με αποτέλεσμα να εγγυάται η σύγκλιση των παραμέτρων στις ιδανικές τιμές (m συχνότητες αρκούν για προσδιορισμό 2m παραμέτρων - PE Condition). Από την ανάλυση του συστήματος με πειραματικές μετρήσεις βρέθηκε πως \(a = -2.59, b = 0.418\). Από τις εξισώσεις των Matching Conditions έχουμε:
    \begin{align*}
        \left.
        \begin{matrix}
            a + bk_x = a_m \\
            bk_r = b_m
        \end{matrix}
        \right\}
        \Rightarrow
        \left.
        \begin{matrix}
            b = \frac{b_m}{k_r} = \frac{0.9}{13}=0.07 \\
            a = \frac{a_m}{bk_m} = \frac{0.9}{0.07 \cdot (-3)}
        \end{matrix}
        \right\}
        \Rightarrow
        \begin{matrix}
            a = -4.3 \\
            b = 0.07 
        \end{matrix}
    \end{align*}
    Παρατηρείται ότι ο ελεγκτής προσεγγίζει πράγματι τις αληθινές τιμές του συστήματος (παίρνοντας υπόψιν το σφάλμα κατά τις πειραματικές μετρήσεις).
    
    \subsection{Ολίσθηση παραμέτρων (drift)}
    % Αφήνοντας τον ελεγκτή να τρέξει για αρκετή ώρα παρατηρείται drift των παραμέτρων:
    Αφήνοντας τον ελεγκτή να τρέξει για αρκετή ώρα παρατηρείται drift των παραμέτρων, όπως φαίνεται στο Διάγραμμα 4.3.1. Αυτό εξηγείται από την θεωρία καθώς η συνάρτηση Lyapunov δημιουργεί μία λωρίδα με φραγμένο σφάλμα παρακολούθησης, αλλά μη φραγμένες εκτιμήσεις παραμέτρων. Έτσι, λόγω θορύβου και μη παραμετρικών αβεβαιοτήτων (όπως αβεβαιότητα μοντέλου) οι εκτιμήσεις των παραμέτρων κάνουν drift προς το άπειρο, ενώ το σφάλμα μένει μικρό. Οι εκτιμήσεις μεγαλώνουν τόσο, ώσπου σε κάποια στιγμή, οι τιμές είναι τόσο μεγάλες που προκαλούν αριθμητική αστάθεια και "έκρηξη" του συστήματος. Αυτό το πρόβλημα διορθώνεται με τροποποιήσεις του νόμου προσαρμογής όπου θα αναλυθεί παρακάτω. % ήταν μετά την εικόνα, αλλά τα ένωσα γιατί υπήρχε τεράστιο κενό
    
    \begin{figure}[h]
        \centering
        \includegraphics[width=0.64\linewidth]{7-mrac_drift_g_01.png}
        \caption{\textbf{Διάγραμμα 4.3.1} Διάγραμμα ολίσθησης παραμέτρων.}
    \end{figure}
    
    % ----------------------------------------------------- ADI ----------------------------------------------------- %
    \section{Εργαστήριο I.4: Έλεγχος ταχύτητας με τη μέθοδο Adaptive Dynamic Inversion (βαθμωτή περίπτωση)}
    Το γραμμικό μοντέλο αναφοράς βάσει του οποίου θα σχεδιαστεί ο ADI ελεγκτής είναι
    
    \begin{equation*}
        \begin{aligned}
            x_m &= -\alpha_m\,x + b_m\,r \\
                &= -0.9\,x + 0.9\,r
        \end{aligned}
    \end{equation*}
    
    \hfill\break
    \hfill\break
    \hfill\break
    \hfill\break
    Επιπλέον, οι νόμοι ελέγχου δίνονται από τις σχέσεις
    
    \begin{equation*}
        \begin{aligned}
            &\dot{\hat{\alpha}} = \gamma_{\alpha}\,\texttt{x}\,\texttt{e} \\
            &\dot{\hat{b}} = \gamma_b\,\texttt{u}\,\texttt{e}
        \end{aligned}
    \end{equation*}
    
    \hfill\break
    ενώ οι τιμές των \(\hat{a},\,\,\hat{b}\) υπολογίζονται με την απλή μέθοδο ολοκλήρωσης του Euler. Αντίστοιχα, ο νόμος ελέγχου είναι
    
    \begin{equation*}
        u = \frac{1}{\hat{b}} [(\alpha_m - \hat{a})x + b_m r]
    \end{equation*}
    
    \subsection{Σήμα αναφοράς: βηματική συνάρτηση}
    Οι εκτιμήσεις των παραμέτρων της προσομοίωσης είναι \(\hat{a} = 2,\,\,\hat{b} = 0.1\).
    
    \begin{figure}[H]
        \centering
        \begin{subfigure}{0.45\textwidth}
            \centering
            \includegraphics[width=\textwidth]{1-adi_step_tracking.png}
            \caption{\textbf{Διάγραμμα 5.1.1} Παρακολούθηση του σήματος αναφοράς \(r(t) = 18\).}
        \end{subfigure}
        \,\,\,\,\,\,
        \begin{subfigure}{0.45\textwidth}
            \centering
            \includegraphics[width=\textwidth]{2-adi_step_drift.png}
            \caption{\textbf{Διάγραμμα 5.1.2} Αποτελέσματα προσομοίωσης με drift στην παράμετρο \(\hat{a}\).}
        \end{subfigure}
    \end{figure}
    
    \hfill\break
    Σύμφωνα με το Διάγραμμα 5.1.1 η παρακολούθηση του σήματος αναφοράς γίνεται ικανοποιητικά, με ένα σφάλμα περίπου 1-2 στροφές το δευτερόλεπτο. Αυτό συμβαίνει γιατί οι εκτιμήσεις των παραμέτρων έχουν "συγκλίνει" σε τιμές που επιτρέπουν στον νόμο ελέγχου να εκμηδενίσει το σφάλμα. Ωστόσο, αξίζει να σημειώσουμε πως η τιμή του \(\gamma_b\) είναι πολύ μικρή\footnote{Η επιλογή της τιμής \(\gamma_b = 10^{-6}\) αποτελεί έναν συμβιβασμό (trade-off) μεταξύ ταχύτητας προσαρμογής της παραμέτρου και ευστάθειας του συστήματος. Παρόλο που περιορίζει τον ελεγκτή, καθιστώντας τον λιγότερο αποδοτικό σε απότομες μεταβολές, κρίνεται απαραίτητη για την αποφυγή του drift της παραμέτρου. Παρακάτω θα μελετήσουμε το σ-modification και e-modification, που βελτιώνουν την παρακολούθηση.} έτσι ώστε το \(\hat{b}\) να μην περάσει από το 0, αφού αλλάζει με πολύ μικρό ρυθμό και μένει κοντά στην αρχική του τιμή με το ίδιο πρόσημο.
    
    \hfil\break
    Σημαντικό είναι και το γεγονός ότι λόγω του φτωχού συχνοτικού περιεχομένου της βηματικής συνάρτησης οι εκτιμήσεις των παραμέτρων μεταβάλλονται αργά με την πάροδο του χρόνου και τελικά η παράμετρος \(\hat{a}\) αποκλίνει (drift), κάτι που γίνεται εμφανές στο Διάγραμμα 5.1.2. Ουσιαστικά ο προσαρμοστικός νόμος προσπαθεί να αντισταθμίσει μικρά σφάλματα αλλάζοντας συνεχώς την τιμή του \(\hat{a}\) χωρίς αυτό να είναι απαραίτητο, οδηγώντας την παράμετρο σε αστάθεια.
    
    \hfil\break
    Τέλος, ας παρατηρήσουμε τη συμπεριφορά του συστήματος κατά την αλλαγή του προσήμου της εκτίμησης των παραμέτρων. Τα αποτελέσματα φαίνονται στο Διάγραμμα 5.1.3. Όταν η παράμετρος \(\hat{b}\) πλησιάσει το 0, πριν πάει να γίνει αρνητική, ο έλεγχος \(u\) απειρίζεται και το σύστημα καταρρέει. Αν από την άλλη καταφέρει να αλλάξει πρόσημο χωρίς να μηδενιστεί (λόγω της διακριτοποίησης των μετρήσεων), ο ελεγκτής θεωρεί ότι το σύστημα έχει πλέον αντίθετη δυναμική από την πραγματική, με αποτέλεσμα να οδηγεί σε μηδενισμό των στροφών. Σε κάθε περίπτωση, η αλλαγή προσήμου του \(\hat{b}\) είναι καταστροφική.
    
    \begin{figure}[h]
        \centering
        \includegraphics[width=0.64\textwidth]{3-adi_step_cross0.png}
        \caption{\textbf{Διάγραμμα 5.1.3} Αποτελέσματα προσομοίωσης για την περίπτωση όπου η παράμετρος \(\hat{b}\) περνά το 0.}
    \end{figure}
    
    \subsection{Σήμα αναφοράς: ημιτονοειδής συνάρτηση}
    Οι εκτιμήσεις των παραμέτρων της προσομοίωσης παραμένουν ίδιες με πριν, ενώ το σήμα αναφοράς είναι \(r(t) = 16 + 4 \sin(0.5 t)\). Σε αντίθεση με πριν, τώρα ικανοποιείται η συνθήκη Persistency of Excitation (PE), καθώς το σήμα εισόδου είναι πλούσιο σε συχνοτικό περιεχόμενο, με αποτέλεσμα οι εκτιμήσεις των παραμέτρων να συγκλίνουν στις πραγματικές τιμές τους. Όπως μπορούμε να παρατηρήσουμε στο Διάγραμμα 5.2.1, η τιμή του \(\hat{a}\) φαίνεται να κάνει ταλάντωση. Αυτό συμβαίνει γιατί το \(x\) εναλλάσσεται συνεχώς, με αποτέλεσμα ο ρυθμός μεταβολής του \(\dot{\hat{a}} = \gamma_a x e\) να αλλάζει και αυτός, σε αντίθεση με τον ρυθμό μεταβολής του \(\dot{\hat{b}} = \gamma_b u e\), που εξαρτάται από το σφάλμα και τον νόμο ελέγχου.
    
    \begin{figure}[h]
        \centering
        \includegraphics[width=0.64\textwidth]{7-adi-PE.png}
        \caption{\textbf{Διάγραμμα 5.2.1} Παρακολούθηση του ημιτονοειδούς σήματος αναφοράς και έλεγχος των PE συνθηκών.}
    \end{figure}
    
    % ------------------------------------------------ MODIFICATIONS ------------------------------------------------ %
    \section{Εργαστήριο I.5: dead-zone, σ-modification, e-modification για τις μεθόδους MRAC και ADI}
    Οι μέθοδοι που θα αναλυθούν παρακάτω είναι παραλλαγές του νόμου προσαρμογής του MRAC και ADI έτσι ώστε να αλλάξει η μορφή της συνάρτησης Lyapunov από n-διάστατη λωρίδα σε n-διάστατο κλειστό σχήμα. Αυτό έχει ως αποτέλεσμα την ευρωστία του ελεγκτή και την αποφυγή του drift των παραμέτρων.
    \subsection{Ανάλυση dead-zone}
    \begin{figure}[H]
        \centering
        \begin{subfigure}{0.45\textwidth}
            \centering
            \includegraphics[width=\textwidth]{13-mrac_dead_zone_step.png}
            \caption{\textbf{Διάγραμμα 6.1.1} Παρακολούθηση του σήματος αναφοράς \(r(t) = 18\) με τον ελεγκτή MRAC με νόμο προσαρμογής dead-zone.}
        \end{subfigure}
        \,\,\,\,\,\,
        \begin{subfigure}{0.45\textwidth}
            \centering
            \includegraphics[width=\textwidth]{4-adi_dead_zone_step.png}
            \caption{\textbf{Διάγραμμα 6.1.2} Παρακολούθηση του σήματος αναφοράς \(r(t) = 18\) με τον ελεγκτή ADI με νόμο προσαρμογής dead-zone.}
        \end{subfigure}
    \end{figure}
    
    \hfill\break
    Με τη μέθοδο αυτή εξασφαλίζουμε πως αν το σφάλμα πέσει κάτω από ένα όριο \(\epsilon\), τότε δεν υπάρχει λόγος να μεταβάλλονται οι παράμετροι \(k_x,\,\,k_r\) (ή \(\hat{a},\,\,\hat{b}\) αντίστοιχα), επομένως ο νόμος ελέγχου παραμένει σταθερός. Στην αντίθετη περίπτωση ανανεώνονται κανονικά οι τιμές των παραμέτρων. Και στις δύο περιπτώσεις γίνεται καλή παρακολούθηση του σήματος αναφοράς, με τις παραμέτρους να συγκλίνουν σε διαφορετικές τιμές. Να σημειώσουμε πως σε κάθε περίπτωση δε παρατηρείται drift σε καμία παράμετρο και αυτό οφείλεται ακριβώς στην εφαρμογή της μεθόδου dead-zone. Το αρνητικό της μεθόδου αυτής είναι το chattering γύρων από τ
    
    \subsection{Ανάλυση σ-modification}
    Η προηγούμενη μέθοδος που χρησιμοποιήσαμε για να διατηρήσουμε τις παραμέτρους εντός καθορισμένων ορίων δεν είναι και ο καλύτερος τρόπος για τη βελτίωση του ελεγκτή μας. Για τον λόγο αυτό χρησιμοποιούμε το σ-modification.
    
    \begin{figure}[H]%[h!]
        \centering
        \begin{subfigure}{0.45\textwidth}
            \centering
            \includegraphics[width=\textwidth]{9-mrac_sigma05_step.png}
            \caption{\textbf{Διάγραμμα 6.2.1} Παρακολούθηση του σήματος αναφοράς \(r(t) = 18\) με τον ελεγκτή MRAC με σ-modification.}
        \end{subfigure}
        \,\,\,\,\,\,
        \begin{subfigure}{0.45\textwidth}
            \centering
            \includegraphics[width=\textwidth]{5-adi_sigma_mod_step.png}
            \caption{\textbf{Διάγραμμα 6.2.2} Παρακολούθηση του σήματος αναφοράς \(r(t) = 18\) με τον ελεγκτή ADI με σ-modification.}
        \end{subfigure}
    \end{figure}
    
    \hfill\break
    Βάσει αυτής της μεθόδου επεμβαίνουμε άμεσα στους ρυθμούς μεταβολής των παραμέτρων ως εξής:
    
    \begin{equation*}
        \begin{aligned}
            &\dot{\hat{k}}_x = -\gamma_x\,(\texttt{x}\,\texttt{e} + \sigma_x\,\hat{k}_x) \\
            &\dot{\hat{k}}_r = -\gamma_r\,(\texttt{r}\,\texttt{e} + \sigma_r\,\hat{k}_r)
        \end{aligned}
    \end{equation*}
    
    \hfill\break
    και
    
    \begin{equation*}
        \begin{aligned}
            &\dot{\hat{\alpha}} = \gamma_{\alpha}\,(\texttt{x}\,\texttt{e} - \sigma_{\alpha}\,\hat{\alpha}) \\
            &\dot{\hat{b}} = \gamma_b\,(\texttt{r}\,\texttt{e} - \sigma_b\,\hat{b})
        \end{aligned}
    \end{equation*}
    
    \hfill\break
    Μπορούμε να παρατηρήσουμε πως όταν οι ρυθμοί μεταβολής \(\dot{\hat{k}}_x\) και \(\dot{\hat{\alpha}}\) εξαρτώνται μόνο από το σφάλμα. Αν υπάρχει έστω και ελάχιστο σφάλμα λόγω θορύβου, οι τιμές των παραμέτρων συνεχώς θα αυξάνονται (ή θα μειώνονται). Με τους όρους \(\sigma\,\hat{p}\), όπου \(p\) η εκάστοτε παράμετρος, αυτό δεν υφίσταται και οι τιμές των παραμέτρων αποσβένουν. Παρατηρούμε ότι και στις δύο περιπτώσεις δεν εμφανίζεται ολίσθηση των παραμέτρων και η μέθοδος είναι αποτελεσματική.
    
    \subsection{Ανάλυση e-modification}
    Μια διαφορετική προσέγγιση είναι το e-modification, σύμφωνα με το οποίο επεμβαίνουμε ξανά στους ρυθμούς μεταβολής των παραμέτρων ως εξής:
    
    \begin{equation*}
        \begin{aligned}
            &\dot{\hat{k}}_x = -\gamma_x\,(\texttt{x}\,\texttt{e} + \sigma_x\,|\texttt{e}|\,\hat{k}_x) \\
            &\dot{\hat{k}}_r = -\gamma_r\,(\texttt{r}\,\texttt{e} + \sigma_r\,|\texttt{e}|\,\hat{k}_r)
        \end{aligned}
    \end{equation*}
    
    \hfill\break
    και
    
    \begin{equation*}
        \begin{aligned}
            &\dot{\hat{\alpha}} = \gamma_{\alpha}\,(\texttt{x}\,\texttt{e} - \sigma_{\alpha}\,|\texttt{e}|\,\hat{\alpha}) \\
            &\dot{\hat{b}} = \gamma_b\,(\texttt{r}\,\texttt{e} - \sigma_b\,|\texttt{e}|\,\hat{b})
        \end{aligned}
    \end{equation*}
    
    \begin{figure}[H]%[h]
        \centering
        \begin{subfigure}{0.45\textwidth}
            \centering
            \includegraphics[width=\textwidth]{11-mrac_epsilon_step.png}
            \caption{\textbf{Διάγραμμα 6.3.1} Παρακολούθηση του σήματος αναφοράς \(r(t) = 18\) με τον ελεγκτή MRAC.}
        \end{subfigure}
        \,\,\,\,\,\,
        \begin{subfigure}{0.45\textwidth}
            \centering
            \includegraphics[width=\textwidth]{6-adi_epsilon_mod_step.png}
            \caption{\textbf{Διάγραμμα 6.3.2} Παρακολούθηση του σήματος αναφοράς \(r(t) = 18\) με τον ελεγκτή ADI.}
        \end{subfigure}
    \end{figure}
    
    \hfill\break
    Εδώ η απόσβεση στις παραμέτρους υφίσταται μόνο όταν υπάρχει σφάλμα, με αποτέλεσμα να προσφέρει μεγαλύτερη ακρίβεια στην παρακολούθηση. Αυτό σημαίνει ότι αν το σύστημα καταφέρει να παρακολουθήσει τέλεια το σήμα αναφοράς, έχοντας δηλαδή μηδενικό σφάλμα, η απόσβεση του ρυθμού μεταβολής των παραμέτρων σταματά και οι τιμές των παραμέτρων δεν αποσβένονται, κάτι το οποίο δε συμβαίνει στο σ-modification. Αν συγκρίνουμε τα Διαγράμματα 6.2.1, 6.2.2 και 6.3.1, 6.3.2 θα διαπιστώσουμε ότι η παρακολούθηση με τον MRAC είναι πολύ καλύτερη με e-modification, ενώ στον ADI η συμπεριφορά είναι αρκετά παρόμοια.
    
    %\subsection{Ανάλυση projection}
    %(...)
    
    \section{Εργαστήριο I.6: Έλεγχος γωνίας και ταχύτητας με τη μέθοδο MRAC (διανυσματική περίπτωση)}
    Το σύστημά μας σε αυτή την περίπτωση διαμορφώνεται ως εξής:
    
    \begin{equation*}
        \begin{bmatrix}
            \dot{\theta} \\
            \dot{\omega}
        \end{bmatrix} = 
        A\,
        \begin{bmatrix}
            \theta \\
            \omega
        \end{bmatrix} +
        \begin{bmatrix}
            0 \\
            1
        \end{bmatrix}u
    \end{equation*}
    
    \noindent
    με τον πίνακα \(A \in \mathbb{R}^{2\times2}\) να είναι άγνωστος. Επιπλέον, σε αντίθεση με το μοντέλο μίας μεταβλητής, εδώ θεωρούμε πως η είσοδος επιδράει στην γωνιακή επιτάχυνση και όχι στην γωνιακή ταχύτητα. Τέλος, αξίζει να σημειωθεί πως για $B = [1 \quad 0]^T$ ο προασρομστικός ελεγκτής δεν δούλευε. Το γραμμικό μοντέλο αναφοράς που επιλέξαμε είναι:

    \begin{equation*}
        \begin{bmatrix}
            \dot{\theta}_m \\
            \dot{\omega}_m
        \end{bmatrix} = 
        A_m\,
        \begin{bmatrix}
            \theta_m \\
            \omega_m
        \end{bmatrix} +
        B_m\,r
        = 
        \begin{bmatrix}
            0 & 0.5 \\
            -2 & -1
        \end{bmatrix}
        \begin{bmatrix}
            \theta_m \\
            \omega_m
        \end{bmatrix} +
        \begin{bmatrix}
            0 \\
            2
        \end{bmatrix}r
    \end{equation*}

    \hfill\break
    Αυτό που αλλάζει με τον MRAC της βαθμωτής περίπτωσης είναι οι νόμοι προσαρμογής, που έχουν ως εξής:

    \begin{equation*}
        \begin{aligned}
            &\dot{\hat{K}}_x = -\Gamma_x\,x\,e^T\,P\,B\,sgn(\Lambda) \\
            &\dot{\hat{K}}_r = -\Gamma_r\,r\,e^T\,P\,B\,sgn(\Lambda)
        \end{aligned}
    \end{equation*}

    \hfill\break
    όπου

    \hfill\break
    \begin{math}
        \Gamma_x = \begin{bmatrix}0.001 & 0 \\ 0 & 0.001\end{bmatrix},\,\,\,\Gamma_r = 0.05 \\[10pt]
        x = \begin{bmatrix}\theta & \omega\end{bmatrix}^T,\,\,\,
        e = \begin{bmatrix}e_{\theta} & e_{\omega}\end{bmatrix}^T,\,\,\,
        r \in \mathbb{R} \\[10pt]
        P = \begin{bmatrix}3 & 0.25 \\ 0.25 & 0.625\end{bmatrix} \\[10pt]
    \end{math}

    \noindent
    Ο πίνακας P προκύπτει από την εξίσωση Lyapunov \(P\,A_m+A_m^T\,P = -Q\). Έτσι, ο νόμος ελέγχου u δαμορφώνεται ως εξής:
    \begin{math}
        u = \dot{\hat{K}}_x \cdot x + \dot{\hat{K}}_r \cdot r
    \end{math}

    \subsection{Σήμα αναφοράς: βηματική συνάρτηση}
    Στο Διάγραμμα 7.2.1 μπορούμε να παρατηρήσουμε ότι ο ελεγκτής κάνει πολύ καλή παρακολούθηση των καταστάσεων του συστήματος, δηλαδή της γωνίας και της γωνιακής ταχύτητας αντίστοιχα. 
    
    \begin{figure}[h]
        \centering
        \includegraphics[width=0.64\textwidth]{1-N_MRAC_simple.png}
        \caption{\textbf{Διάγραμμα 7.1.1} Παρακολούθηση σταθερής γωνίας και ταχύτητας αναφοράς: \(90\degree\) και 0 RPM αντίστοιχα.}
    \end{figure}

    \noindent
    Χρησιμοποιώντας τώρα σ-modification έχουμε:

    \begin{figure}[h]
        \centering
        \includegraphics[width=0.64\textwidth]{2-N_MRAC_sigma_mod_step.png}
        \caption{\textbf{Διάγραμμα 7.1.2} Παρακολούθηση γωνίας και ταχύτητας αναφοράς με χρήση σ-modification. Η γωνία ξεκινάει από τις \(90\degree\), ανεβαίνει στις \(180\degree\) και έπειτα πέφτει στις \(120\degree\).}
    \end{figure}

    \noindent
    Εδώ οι νόμοι προσαρμογής γίνονταi αντίστοιχα

    \begin{equation*}
        \begin{aligned}
            &\dot{\hat{K}}_x = -\Gamma_x\,(x\,e^T\,P\,B + \sigma_x\,\hat{K}_x)\,sgn(\Lambda) \\
            &\dot{\hat{K}}_r = -\Gamma_r\,(r\,e^T\,P\,B+ \sigma_r\,\hat{K}_r)\,sgn(\Lambda)
        \end{aligned}
    \end{equation*}

    \hfill\break
    Στο Διάγραμμα 7.1.2 φαίνεται ότι ο ελεγκτής καταφέρνει να ακολουθήσει τα σήματα αναφοράς που ανατίθενται διαδοχικά, έχοντας βέβαια αποκλίσεις στην αρχή, που μειώνονται με τον χρόνο. Οι παράμετροι προσαρμογής από την άλλη παρουσιάζουν έντονες μεταβολές στην αρχή, αλλά τελικά συγκλίνουν σε πιο σταθερές τιμές, επιτυγχάνοντας τον στόχο.
    
    \subsection{Σήμα αναφοράς: ημιτονοειδής συνάρτηση}
    Αντίστοιχα συμπεράσματα μπορούμε να έχουμε και για παρακολούθηση ημιτονοειδoύς σήματος.

    \begin{figure}[H]
        \centering
        \includegraphics[width=0.64\textwidth]{3-N_MRAC_sigma_mod_sin.png}
        \caption{\textbf{Διάγραμμα 7.2.1} Παρακολούθηση του σήματος αναφοράς \(r(t) = 90 + 90 \cdot \sin(0.5\,t)\) με χρήση σ-modification.}
    \end{figure}

    \subsection{Ολίσθηση παραμέτρων και PE συνθήκες}
    Από τις παραπάνω προσομοιώσεις και αφήνοντας τις προσομοιώσεις να τρέξουν για αρκετό χρόνο δεν παρατηρήθηκε drift. Αντίθετα, μετά από κάποιο χρόνο όλες οι παράμετροι σύγκλιναν σε σταθερές τιμές, ενώ το σύστημα παρέμενε ευσταθές. Ειδικότερα, λόγω του ότι ο πίνακας Α του συστήματος είναι άγνωστος δε μπορούν να βρεθούν ιδανικές τιμές των \(K_x\), \(K_r\) για να συγκρίνουμε τα αποτελέσματά μας και να δούμε αν οι τιμές των παραμέτρων κατά την προσομοίωση συγκλίνουν σε αυτές.

    \subsection{Το πρόσημο του \(\Lambda\)}
    Με δοκιμή για θετικό και αρνητικό \(\Lambda\) καταλήξαμε ότι το πρόσημό του είναι αρνητικό για να είναι το σύστημα ευσταθές και να παρακολουθεί τα αντίστοιχα σήματα αναφοράς.
\end{document}
 